%% ma: L10-B12-0005
%% lop: 10
%% chuong: 5
%% bai: 12
%% dang: 10.12.7
%% loai: TL
%% muc: TH
%% dap_an: "-"
%% nguon: "chuyen-de-cac-so-dac-trung-cua-mau-so-lieu-khong-ghep-nhom-toan-10-chuong-trinh-moi.pdf (GV. Phan Nhật Linh) (P01), Bài 01 Bài tập tự luận 5"
%% kien_thuc: "Lớp 10 Bài 12 (số gần đúng, sai số tuyệt đối, độ chính xác, sai số tương đối, quy tròn số)"
%% trang_thai: chinh-thuc
%% ghi_chu: "Tự duyệt theo yêu cầu giáo viên 22/10/2026."
\begin{ex}
Cho tam giác $ABC$ có độ dài ba cạnh đo được là $a=12\text{ cm}\pm0{,}2\text{ cm}$; $b=10{,}2\text{ cm}\pm0{,}2\text{ cm}$; $c=8\text{ cm}\pm0{,}1\text{ cm}$. Tính chu vi $P$ của tam giác và đánh giá sai số tuyệt đối của chu vi qua phép đo.
\loigiai{Viết $a=12+d_1$, $b=10{,}2+d_2$, $c=8+d_3$ với $|d_1|\le0{,}2$, $|d_2|\le0{,}2$, $|d_3|\le0{,}1$.

$P=30{,}2+d_1+d_2+d_3$ và $|d_1+d_2+d_3|\le0{,}5$.

Vậy $P=30{,}2\text{ cm}\pm0{,}5\text{ cm}$, sai số tuyệt đối $\Delta_P\le0{,}5$ cm.}
\end{ex}
